\section{Introduction}
\label{sec:intro}

Grokking, where generalisation arrives long after a training set is fit, is useful here as an
instrument rather than as a subject. A small transformer trained to convergence on a closed
algebraic task can be reverse-engineered completely, and the resulting circuit is a statement
about what the architecture does with the structure it is given
\citep{power2022grokking,nanda2023progress}.

Modular multiplication is the harder case, and a change of basis is the source. For
$a + b \bmod n$ the network represents inputs by characters of $\mathbb{Z}/n\mathbb{Z}$ and
adds phases. For $a \cdot b \bmod n$ the same circuit exists but lives in the multiplicative
character basis, reached by discrete log, and a model that is sharply sparse there looks dense
in raw residue coordinates \citep{nguyen2026dlogclock}. The basis is not a presentational
choice: a statistic read in the wrong one inverts the conclusion.

That account needs the units to form a group, and for a general modulus most of the
multiplication table is not units. \citet{chen2026multiplication} took the step past it: they
model $(\mathbb{Z}/n\mathbb{Z}, \cdot)$ as a monoid, show that character-style computation
localises to disjoint algebraic strata identified with Green's $\mathcal{J}$-classes, and
demonstrate class-sensitive routing. That framing is theirs. Independent derivation is
not claimed and would not be a defence: \citet{chen2026multiplication} published the stratified
framing before this work reached it, and they are cited as its origin. Their result covers four
moduli, all square-free, and they state its two limitations themselves: the evidence is
\emph{only correlational}, and non-square-free moduli introduce non-regular
$\mathcal{J}$-classes containing nilpotent elements, which they describe as breaking the
reduction to local group characters. This paper answers both. It is the non-square-free
completion of their programme, and it replaces correlation with an ablation: the first
of the two instruments they ask for, not the second (Sections~\ref{sec:causal} and
\ref{sec:strata}).

The reduction does not break. Writing $x = d \cdot u$ makes a
$\mathcal{J}$-class a torsor under the \emph{local} units, so a stratum is multiplication
in a smaller ring on that stratum's own local group $G_{m}$. Nilpotents \emph{collapse} a stratum onto a
smaller local group; they do not destroy the character structure
(Theorem~\ref{thm:stratum}, Section~\ref{sec:strata}). The algebra of $n$
does not change the mechanism. It changes how many strata there are and how they nest: one
for a field, a nested $p$-adic chain for a prime power, a divisor lattice for a square-free
composite.

% STE procedural
\begin{figure}[t]
  \centering
  \includegraphics[width=\textwidth]{../figures/F0_pipeline.pdf}
  \caption[What this paper does.]{\textbf{What this paper does.} Train a one-layer transformer to grok
    $a \cdot b \bmod n$ at each of $23$ moduli. Then read the trained model in the coordinate
    the algebra supplies. The first three steps are standard. The fourth is not. Relabel each
    unit by its exponent tuple. Only then does multiplication of units become addition.
    Theorem~\ref{thm:stratum} gives every stratum $J_{d} \times J_{e}$ its own such
    coordinate, on its own local group $G_{m}$.\par
    Measure and ablate per stratum, in that stratum's coordinate and not in the raw residues
    of $\mathbb{Z}/n\mathbb{Z}$. The figure regenerates from
    \texttt{src/viz/plots.py::pipeline}.}
  \label{fig:pipeline}
\end{figure}
% STE end

\paragraph{Contributions.}
\begin{enumerate}
  \item \emph{The stratum identity} (Theorem~\ref{thm:stratum}), proved and also
    checked by exhaustion, which
    turns ``non-regular classes break the reduction'' into ``non-regular classes carry the
    reduction on a smaller group'', and fixes what to measure without any freedom of choice.
  \item \emph{The mechanism, measured on every stratum}: $32$ strata across six moduli, $157$
    stratum-run measurements and $12$ distinct local groups, with criteria committed before any
    per-stratum number existed. It raises the fraction of the multiplication table accounted
    for from $23.5\%$ to $82.3\%$ at $n = 165$ and from $64.0\%$ to $89.6\%$ at $n = 125$
    (Section~\ref{sec:strata}, five seeds).
  \item \emph{Ablation evidence across $23$ moduli, $14$ of them non-square-free}, including
    all six prime powers (Section~\ref{sec:causal}, five seeds at six moduli, three at the
    rest). It reaches $2^{7}$, nilpotency depth $7$, where the key set is beyond all
    $10{,}000$ draws in $2/2$ runs (improvement $2.77$--$476\times$) in product-character
    coordinates where no discrete log exists.
  \item \emph{One mechanism, not one per stratum}: strata that share a local group select
    exactly the same characters of it, $158$ of $160$ pairs across five seeds, with a
    cell-shuffle control at $0$ of $160$. This is the test that separates a single indexed
    mechanism from a private sub-circuit per block (Section~\ref{sec:strata-sameness}). It
    ships without a failed-run control, and Section~\ref{sec:strata-sameness} says why.
  \item \emph{The negative results, reported as results}. The pre-registered primary
    criterion of our own stratum test does not discriminate a grokked model from a failed one,
    and the verdict rests on the three criteria that do. The early CRT enrichment ramp is not a
    progress measure. Five earlier claims of ours are retracted in
    Appendix~\ref{app:retractions}, with what killed each and the control that would have
    caught it.
\end{enumerate}

Section~\ref{sec:related} relates this work to prior mechanistic accounts.
Section~\ref{sec:setup} defines the task and its strata, and Section~\ref{sec:methods} the
model, the coordinates, the statistics and the ablation protocol. Section~\ref{sec:calib}
validates the pipeline against published results and compares our two implementations.
Section~\ref{sec:clock} reports where models grok and the clock at prime powers,
Section~\ref{sec:causal} the ablation evidence, Section~\ref{sec:strata} the measurement on
every stratum and Section~\ref{sec:crt} the additive CRT structure.
Section~\ref{sec:discussion} discusses the results, Section~\ref{sec:limits} states the
limitations and Section~\ref{sec:conclusion} concludes.
